\documentclass[11pt]{article}

%==================== packages ====================
\usepackage[T1]{fontenc}
\usepackage[protrusion=true,expansion=false]{microtype}
\usepackage[a4paper,margin=1in]{geometry}
\usepackage{amsmath,amssymb,amsthm,mathtools}
\usepackage{xcolor}
\usepackage{enumitem}
\usepackage{titlesec}
\usepackage[hidelinks]{hyperref}

%==================== brand palette ====================
\definecolor{navy}{HTML}{1A1A2E}
\definecolor{sky}{HTML}{3FA9F5}
\definecolor{sage}{HTML}{8BB999}
\definecolor{gold}{HTML}{F3D630}

\hypersetup{
  colorlinks=true,
  linkcolor=navy,
  citecolor=navy,
  urlcolor=sky,
}

\titleformat{\section}{\normalfont\large\bfseries\color{navy}}{\thesection}{0.6em}{}
\titleformat{\subsection}{\normalfont\bfseries\color{navy}}{\thesubsection}{0.6em}{}

%==================== theorem environments ====================
\theoremstyle{definition}
\newtheorem{definition}{Definition}[section]
\newtheorem{notation}[definition]{Notation}
\newtheorem{remark}[definition]{Remark}
\newtheorem{example}[definition]{Example}

\theoremstyle{plain}
\newtheorem{proposition}[definition]{Proposition}
\newtheorem{lemma}[definition]{Lemma}
\newtheorem{corollary}[definition]{Corollary}

%==================== notation macros ====================
\newcommand{\Ob}{\mathsf{O}}            % sort of boundaries
\newcommand{\Ed}{\mathsf{E}}            % sort of elements
\newcommand{\M}{M}                      % comprehension former
\newcommand{\unit}{\eta}                % unit / head
\newcommand{\zero}{\varnothing}         % zero / empty
\newcommand{\guard}[1]{\langle #1\rangle}
\newcommand{\eq}{\mathrel{\doteq}}      % guard equality
\newcommand{\gen}[2]{#1\!\leftarrow\!#2}% generator
\newcommand{\dl}{\partial_{0}}          % left boundary
\newcommand{\dr}{\partial_{1}}          % right boundary
\newcommand{\uu}{\mathsf{u}}            % neutral-element operation
\newcommand{\kk}{\kappa}                % multiplication
\newcommand{\Hom}{\mathsf{H}}           % derived hom-comprehension
\newcommand{\one}{\mathbf{1}}           % terminal sort
\newcommand{\pt}{\ast}                  % its element
\newcommand{\id}{\mathrm{id}}
\newcommand{\co}{\mathbin{\diamond}}    % fibred multiplication
\newcommand{\unitf}{\mathsf{e}}         % fibred neutral element
\newcommand{\compr}[2]{\bigl[\,#1 \ \big|\ #2\,\bigr]}

\setlist[enumerate]{itemsep=2pt,topsep=3pt}

%==================== title ====================
\title{\color{navy}\bfseries Boundaried Monoids via Comprehension\\[2pt]
\normalsize\normalfont\color{black} A purely equational theory}
\author{L.\,G.\,Meredith\\[-1pt]\small F1R3FLY.io}
\date{June 2026}

\begin{document}
\maketitle

\begin{abstract}
\noindent
We give a self-contained, purely equational theory $\mathbf{T}$ of what we call
\emph{boundaried monoids}: a generalisation of monoids in which every element
carries a left and a right boundary drawn from an auxiliary sort, and the
product of two elements is defined exactly when the inner boundaries agree.
The development uses no notion beyond sorts, operations, equations, and
\emph{comprehension}. Its one structural idea is presentational: rather than
treat the product as a partial operation hedged by side conditions, we absorb
its partiality into a \emph{comprehension structure}, so that the product
becomes a total operation valued in a collection sort. The translation between
the partial and the total readings is exactly comprehension in the sense of
Wadler~\cite{wadler}. We axiomatise the comprehension layer (\S\ref{sec:comp}),
state the principal theory over it (\S\ref{sec:T}), give an equivalent
dependently sorted presentation, and prove the two equivalent with comprehension
as the translating device (\S\ref{sec:equiv}). Throughout, the object of study
is an algebra and nothing more.
\end{abstract}

%================================================================
\section{Aim and standing conventions}\label{sec:intro}

This note isolates a single algebraic theory and writes down its equations.
We deliberately confine ourselves to elementary algebraic vocabulary: a
\emph{signature} is a family of sorts together with operations between finite
products of sorts; a \emph{theory} adds equations; a \emph{model} is a
sort-indexed family of carriers interpreting the operations and satisfying the
equations. Beyond this we use only one further device, \emph{comprehension},
which we axiomatise in \S\ref{sec:comp} before it is used.

The structure we axiomatise is a mild generalisation of a monoid. Recall that a
monoid is a set with an associative binary product and a two-sided neutral
element. A \emph{boundaried monoid} relaxes totality of the product: each
element $f$ is assigned a left boundary $\dl f$ and a right boundary $\dr f$ in
an auxiliary sort $\Ob$, and the product $g\cdot f$ is intended to exist
precisely when $\dl g=\dr f$. A monoid is the special case in which $\Ob$ has a
single element, so that every pair is multipliable.

There are two faithful ways to render the partial product. The first keeps the
product partial and indexes it by boundaries, giving a dependently sorted
\emph{equational} theory with totally standard operations
(\S\ref{sec:equiv}). The second keeps a single sort of elements and makes the
product \emph{total} by valuing it in a collection sort: the product of a
non-multipliable pair is the empty collection, and the product of a
multipliable pair is a singleton. The second reading is the one in which the
partiality is handled by comprehension, and it is the reading we take as
primary. The bulk of the note shows that the two readings define the same
theory.

\paragraph{A note on scope.} We make no semantic claims here; interpretations of
$\mathbf{T}$ are deferred to separate work. We also make no use of any structure
richer than the comprehension layer of \S\ref{sec:comp}. The reader will
recognise that layer, under a different packaging, as a familiar gadget; we cite
its provenance~\cite{wadler,moggi} and otherwise do not rely on that
packaging.\footnote{The comprehension structure of
Definition~\ref{def:comp} is an \emph{extension system}: a unit together with a
Kleisli-style extension subject to three equations. We present it in exactly
this elementary form so that the present theory remains an algebra, with no
appeal to functoriality or to ambient morphisms.}

%================================================================
\section{The comprehension layer}\label{sec:comp}

We work over a fixed many-sorted signature that is closed under finite products:
it has a terminal sort $\one$ with unique element $\pt$, binary products
$X\times Y$ with projections and pairing $\langle\,,\rangle$, and the evident
equations. We write $f\colon X\to Y$ for a derived operation (a term in one free
variable). All of this is ordinary equational algebra.

\begin{definition}[Comprehension structure]\label{def:comp}
A \emph{comprehension structure} consists of
\begin{enumerate}[label=(\roman*)]
  \item a sort $\M X$ for every sort $X$ (the sort of \emph{collections} of $X$);
  \item for each sort $X$ a \emph{unit} $\unit_X\colon X\to\M X$;
  \item for each operation $f\colon X\to\M Y$ an \emph{extension}
        $f^{\star}\colon \M X\to\M Y$;
  \item for each sort $X$ a constant $\zero_X\colon\M X$ (the \emph{empty
        collection});
  \item an \emph{equality guard}
        $\guard{-\eq-}\colon \Ob\times\Ob\to\M\one$
        on the distinguished sort $\Ob$,
\end{enumerate}
subject to the equations
\begin{align}
  \unit_X^{\star} &= \id_{\M X}, \tag{C1}\\
  f^{\star}\circ\unit_X &= f, \tag{C2}\\
  (g^{\star}\circ f)^{\star} &= g^{\star}\circ f^{\star}, \tag{C3}
\end{align}
for all $f\colon X\to\M Y$, $g\colon Y\to\M Z$; the \emph{commutativity} law
\begin{equation}
  \compr{(x,y)}{\gen{x}{u},\,\gen{y}{v}}
  \;=\;
  \compr{(x,y)}{\gen{y}{v},\,\gen{x}{u}}
  \tag{C4}
\end{equation}
whenever $u\colon\M X$ and $v\colon\M Y$ share no free variable; the \emph{zero}
laws
\begin{align}
  f^{\star}\circ\zero_X &= \zero_Y, \tag{C5}\\
  \compr{t}{\gen{x}{\zero_X},\,\overline{q}} &= \zero, \tag{C6}
\end{align}
and the \emph{guard} laws
\begin{align}
  \guard{a\eq a} &= \unit_{\one}(\pt), \tag{C7}\\
  \compr{\varphi(a,b)}{\guard{a\eq b},\,\overline{q}}
  &= \compr{\varphi(a,a)}{\guard{a\eq b},\,\overline{q}}. \tag{C8}
\end{align}
The comprehension notation $\compr{\,\cdot\,}{\,\cdot\,}$ used above is defined in
Notation~\ref{not:comp}; equations (C4)--(C8) are to be read after that
expansion.
\end{definition}

\begin{notation}[Comprehension]\label{not:comp}
A \emph{qualifier} is either a \emph{generator} $\gen{x}{u}$ (with
$u\colon\M X$ and $x$ a fresh variable of sort $X$) or a \emph{guard}
$\guard{a\eq b}$ (with $a,b\colon\Ob$). For a term $t$ and a finite list
$\overline{q}=q_1,\dots,q_n$ of qualifiers we define a term
$\compr{t}{\overline{q}}\colon\M Y$ (where $t\colon Y$ in the context bound by
$\overline{q}$) by recursion on the list:
\begin{align*}
  \compr{t}{\,} &\;=\; \unit(t),\\
  \compr{t}{\gen{x}{u},\,\overline{q}}
    &\;=\; \bigl(\lambda x.\,\compr{t}{\overline{q}}\bigr)^{\!\star} u,\\
  \compr{t}{\guard{a\eq b},\,\overline{q}}
    &\;=\; \bigl(\lambda \pt.\,\compr{t}{\overline{q}}\bigr)^{\!\star}\,
           \guard{a\eq b}.
\end{align*}
Thus generators are extension, the head is the unit, and guards filter. We read
$\compr{t}{\overline q}$ as ``collect $t$ as the bound variables range over the
generators, retaining only those bindings that satisfy the guards.''
\end{notation}

The equations have transparent readings. (C1)--(C3) say that collecting a
singleton is identity, that a singleton generator binds by substitution, and
that nested collection may be flattened in either order: they are the laws that
make comprehension well defined. (C4) says independent generators may be
reordered. (C5)--(C6) say that nothing is collected over an empty range, from
either side. (C7) says a reflexive guard is vacuous (it contributes the
singleton $\unit(\pt)$, which passes through); (C8) says that under a guard
$\guard{a\eq b}$ one may rewrite $b$ to $a$. We record the two immediate
consequences we shall need.

\begin{lemma}[Reduction]\label{lem:reduce}
For all suitable terms,
\begin{align}
  \compr{t}{\guard{a\eq a},\,\overline q}
    &= \compr{t}{\overline q}, \tag{R1}\label{eq:R1}\\
  \compr{t}{\gen{x}{\unit(s)},\,\overline q}
    &= \compr{t[s/x]}{\overline q}. \tag{R2}\label{eq:R2}
\end{align}
\end{lemma}
\begin{proof}
(R1): by Notation~\ref{not:comp} the left side is
$(\lambda\pt.\,\compr{t}{\overline q})^{\star}\guard{a\eq a}$, which by (C7)
equals $(\lambda\pt.\,\compr{t}{\overline q})^{\star}(\unit\,\pt)$, and this is
$\compr{t}{\overline q}$ by (C2). (R2): the left side is
$(\lambda x.\,\compr{t}{\overline q})^{\star}(\unit\,s)$, which is
$\compr{t}{\overline q}[s/x]=\compr{t[s/x]}{\overline q}$ by (C2).
\end{proof}

\begin{notation}[Subsort by comprehension]\label{not:subsort}
For a sort $X$ and parallel terms $p,q\colon X\to\Ob$ we write
\[
  \{\,x\colon X \mid p(x)\eq q(x)\,\}
\]
for the carrier whose elements are those $x\colon X$ with $p(x)=q(x)$, with the
inclusion term and the universal property of an equationally defined subsort.
Under the comprehension reading this subsort is the support of
$\compr{x}{\gen{x}{w},\,\guard{p(x)\eq q(x)}}$ for $w$ ranging over $\M X$; the
guard discards the elements failing the equation, by (C6)--(C8). We use the
subsort notation freely below as a convenient name for this support.
\end{notation}

%================================================================
\section{The theory \texorpdfstring{$\mathbf{T}$}{T} of boundaried monoids}
\label{sec:T}

We now state the principal theory over a fixed comprehension structure.

\begin{definition}[Signature]\label{def:sig}
The signature of $\mathbf{T}$ extends that of \S\ref{sec:comp} with
\begin{enumerate}[label=(\roman*)]
  \item two sorts $\Ob$ (\emph{boundaries}) and $\Ed$ (\emph{elements}), with
        $\Ob$ the distinguished sort carrying the guard of
        Definition~\ref{def:comp};
  \item \emph{boundary} operations $\dl,\dr\colon\Ed\to\Ob$;
  \item a \emph{neutral} operation $\uu\colon\Ob\to\Ed$;
  \item a \emph{multiplication}
        $\kk\colon\Ed\times\Ed\to\M\Ed$.
\end{enumerate}
We read $\kk(g,f)$ as ``the collection of products of $g$ after $f$,'' a total
operation whose value is empty when $f,g$ are not multipliable and a singleton
otherwise.
\end{definition}

\begin{definition}[Derived collections]\label{def:derived}
Over the signature of Definition~\ref{def:sig} we name three comprehensions.
\begin{align*}
  \Hom(a,b) &\;:=\; \{\,f\colon\Ed \mid \dl f\eq a,\ \dr f\eq b\,\}
            &&\text{(elements with prescribed boundaries),}\\
  \Ed^{(2)} &\;:=\; \{\,(g,f)\colon\Ed\times\Ed \mid \dl g\eq \dr f\,\}
            &&\text{(multipliable pairs),}\\
  \Ed^{(3)} &\;:=\; \{\,(h,g,f) \mid \dl h\eq\dr g,\ \dl g\eq\dr f\,\}
            &&\text{(multipliable triples).}
\end{align*}
Each is a subsort in the sense of Notation~\ref{not:subsort}; none is a new
primitive. The multiplication $\kk$ is defined on all of $\Ed\times\Ed$, but the
axioms below force it to be supported on $\Ed^{(2)}$.
\end{definition}

\begin{definition}[Axioms of $\mathbf{T}$]\label{def:axioms}
A model of $\mathbf{T}$ is a comprehension structure together with
interpretations of $\Ob,\Ed,\dl,\dr,\uu,\kk$ satisfying the following equations,
for all $a\colon\Ob$ and $f,g,h\colon\Ed$.

\medskip
\noindent\emph{Boundaries of the neutral element.}
\begin{align}
  \dl\bigl(\uu(a)\bigr) &= a, \tag{T1}\\
  \dr\bigl(\uu(a)\bigr) &= a. \tag{T2}
\end{align}

\noindent\emph{Support and determinacy of the product.}
\begin{align}
  \compr{\pt}{\gen{v}{\kk(g,f)}}
    &= \guard{\dl g\eq\dr f}, \tag{T3}\\
  \compr{(v,v')}{\gen{v}{\kk(g,f)},\,\gen{v'}{\kk(g,f)}}
    &= \compr{(v,v)}{\gen{v}{\kk(g,f)}}. \tag{T4}
\end{align}

\noindent\emph{Boundaries of a product.}
\begin{align}
  \compr{\dl v}{\gen{v}{\kk(g,f)}}
    &= \compr{\dl f}{\gen{v}{\kk(g,f)}}, \tag{T5}\\
  \compr{\dr v}{\gen{v}{\kk(g,f)}}
    &= \compr{\dr g}{\gen{v}{\kk(g,f)}}. \tag{T6}
\end{align}

\noindent\emph{Neutrality.}
\begin{align}
  \kk\bigl(\uu(\dr f),\,f\bigr) &= \unit(f), \tag{T7}\\
  \kk\bigl(f,\,\uu(\dl f)\bigr) &= \unit(f). \tag{T8}
\end{align}

\noindent\emph{Associativity.}
\begin{equation}
  \compr{w}{\gen{v}{\kk(g,f)},\,\gen{w}{\kk(h,v)}}
  \;=\;
  \compr{w}{\gen{v}{\kk(h,g)},\,\gen{w}{\kk(v,f)}}. \tag{T9}
\end{equation}
\end{definition}

A few words on why these nine equations say what is intended.

(T3) is the \emph{support} law: mapping the collection $\kk(g,f)$ onto the
terminal sort yields $\guard{\dl g\eq\dr f}$, which is the singleton
$\unit(\pt)$ when $g,f$ are multipliable and is empty otherwise (by (C7) and the
discussion of Notation~\ref{not:subsort}). Hence $\kk(g,f)$ is non-empty exactly
on $\Ed^{(2)}$.

(T4) is \emph{determinacy}: any two elements drawn from $\kk(g,f)$ coincide, so
$\kk(g,f)$ has at most one element. With (T3) this makes $\kk(g,f)$ a singleton
on $\Ed^{(2)}$ and empty off it. Partiality has been wholly absorbed into the
collection sort; $\kk$ itself is total.

(T5)--(T6) transport boundaries across a product: wherever a product exists, its
left boundary is that of $f$ and its right boundary is that of $g$. Stating
these by mapping over $\kk(g,f)$ keeps them equational and vacuously true off
$\Ed^{(2)}$.

(T7)--(T8) are the two neutral laws. The pair on the left of (T7) is multipliable
since $\dl(\uu(\dr f))=\dr f$ by (T1), so the guard in (T3) passes and the law
asserts that the product is the singleton $\unit(f)$; symmetrically for (T8).

(T9) is associativity. Both sides are double comprehensions: the left forms
$g$ after $f$, then $h$ after that; the right forms $h$ after $g$, then that
after $f$. By (T3) each side is non-empty exactly when
$\dl g\eq\dr f$ and $\dl h\eq\dr g$, the common condition defining
$\Ed^{(3)}$, and on that subsort each is a singleton; the equation asserts the
two singletons agree. Note that the boundary bookkeeping that ordinarily
accompanies associativity is here discharged automatically by the guards inside
$\kk$, via (T3) and (T5)--(T6).

\begin{remark}[Monoids as a special case]
If $\Ob\cong\one$ then $\dl,\dr$ are constant, every pair is multipliable, (T3)
makes $\kk(g,f)$ a singleton for all $g,f$, and (T7)--(T9) reduce to the
ordinary unit and associativity laws of a monoid with unit $\uu(\pt)$. Thus
$\mathbf{T}$ specialises to the equational theory of monoids.
\end{remark}

\begin{lemma}[Single-valued product]\label{lem:single}
In any model of $\mathbf{T}$ there is, for each multipliable pair
$(g,f)\in\Ed^{(2)}$, a unique element $g\bullet f\colon\Ed$ with
$\kk(g,f)=\unit(g\bullet f)$, and it satisfies $\dl(g\bullet f)=\dl f$,
$\dr(g\bullet f)=\dr g$.
\end{lemma}
\begin{proof}
On $\Ed^{(2)}$ the guard $\guard{\dl g\eq\dr f}$ equals $\unit(\pt)$ by (C7), so
(T3) gives $\compr{\pt}{\gen{v}{\kk(g,f)}}=\unit(\pt)$: the collection
$\kk(g,f)$ is non-empty. By (T4) it is a subsingleton. A non-empty subsingleton
is a singleton $\unit(g\bullet f)$ for a unique $g\bullet f$; uniqueness is by
(C2). The boundary identities follow from (T5)--(T6) and (R2) applied to
$\kk(g,f)=\unit(g\bullet f)$.
\end{proof}

%================================================================
\section{The fibred presentation, and equivalence}\label{sec:equiv}

We now give the second, dependently sorted reading and show it defines the same
theory. Here the product is honestly total and single-valued, and no zero or
guard is needed; the price is a sort that depends on two parameters.

\begin{definition}[The fibred theory $\mathbf{T}^{\flat}$]\label{def:flat}
The theory $\mathbf{T}^{\flat}$ has
\begin{enumerate}[label=(\roman*)]
  \item a sort $\Ob$, and for each $a,b\colon\Ob$ a sort $\Hom(a,b)$;
  \item neutral elements $\unitf_a\colon\Hom(a,a)$ for each $a\colon\Ob$;
  \item products
        $\co_{a,b,c}\colon\Hom(b,c)\times\Hom(a,b)\to\Hom(a,c)$
        for each $a,b,c\colon\Ob$;
\end{enumerate}
subject to the equations, for $f\colon\Hom(a,b)$, $g\colon\Hom(b,c)$,
$h\colon\Hom(c,d)$,
\begin{align}
  \unitf_b \co f &= f, \tag{F1}\\
  f \co \unitf_a &= f, \tag{F2}\\
  h \co (g\co f) &= (h\co g)\co f. \tag{F3}
\end{align}
This is an ordinary equational theory, with total single-valued operations and
three equations; it uses no comprehension structure.
\end{definition}

\begin{proposition}[Equivalence]\label{prop:equiv}
$\mathbf{T}$ and $\mathbf{T}^{\flat}$ are definitionally equivalent: there are
mutually inverse interpretations of each theory in the other, definable using
the comprehension structure, that are the identity on $\Ob$.
\end{proposition}

\begin{proof}[Proof sketch]
\emph{From $\mathbf{T}$ to $\mathbf{T}^{\flat}$.} Keep $\Ob$. Define
$\Hom(a,b)$ by the comprehension of Definition~\ref{def:derived}. Put
$\unitf_a:=\uu(a)$, which lies in $\Hom(a,a)$ by (T1)--(T2). For
$g\in\Hom(b,c)$ and $f\in\Hom(a,b)$ the pair $(g,f)$ is multipliable, so
Lemma~\ref{lem:single} supplies a unique $g\bullet f$ with
$\kk(g,f)=\unit(g\bullet f)$, lying in $\Hom(a,c)$ by the boundary identities;
set $g\co f:=g\bullet f$. Equations (F1)--(F2) are (T7)--(T8) with the unit
stripped by (C2). For (F3): apply (R2) twice inside each side of (T9),
replacing $\kk(g,f)$ by $\unit(g\bullet f)$ and then $\kk(h,g\bullet f)$ by
$\unit(h\bullet(g\bullet f))$ on the left, and symmetrically on the right; (T9)
then reads $\unit\bigl(h\bullet(g\bullet f)\bigr)=\unit\bigl((h\bullet g)\bullet
f\bigr)$, whence (F3) by (C2).

\emph{From $\mathbf{T}^{\flat}$ to $\mathbf{T}$.} Keep $\Ob$. Define the sort of
elements as the comprehension (disjoint collection)
\[
  \Ed \;:=\; \{\,(a,b,f) \mid a,b\colon\Ob,\ f\colon\Hom(a,b)\,\},
\]
with $\dl(a,b,f):=a$, $\dr(a,b,f):=b$, and $\uu(a):=(a,a,\unitf_a)$. Define the
multiplication as the guarded singleton
\[
  \kk\bigl((b',c,g),(a,b,f)\bigr)
   \;:=\;
   \compr{\,(a,c,\;g\co f)\,}{\guard{b'\eq b}},
\]
which by (C7) is $\unit(a,c,g\co f)$ when $b'=b$ and is $\zero$ otherwise (the
$g\co f$ on the right is well typed precisely under the guard, after rewriting
$b'$ to $b$ by (C8)). One checks the axioms directly: (T3) and (T4) hold because
the value is a guarded singleton; (T5)--(T6) because the boundaries of
$(a,c,g\co f)$ are $a$ and $c$; (T7)--(T8) from (F1)--(F2); and (T9) from (F3),
the two guards on each side combining to the single condition defining
$\Ed^{(3)}$.

\emph{Mutual inverse.} Starting from $\mathbf{T}$, forming $\mathbf{T}^{\flat}$
and then $\mathbf{T}$ again returns $\Ed$ up to the evident identification
$f\mapsto(\dl f,\dr f,f)$, under which $\dl,\dr,\uu,\kk$ are preserved on the
nose by Lemma~\ref{lem:single}. Starting from $\mathbf{T}^{\flat}$, the round
trip returns $\Hom(a,b)$ as the comprehension
$\{(a',b',f)\mid a'\eq a, b'\eq b\}$, naturally identified with $\Hom(a,b)$ by
(C8). Both identifications respect all operations, so the interpretations are
mutually inverse. The only non-equational devices used in either direction are
comprehensions, which is the asserted translating role of comprehension.
\end{proof}

\begin{corollary}
$\mathbf{T}$ is conservative over the comprehension layer of
\S\ref{sec:comp}: any equation between terms not mentioning $\dl,\dr,\uu,\kk$
that holds in $\mathbf{T}$ already holds in the comprehension layer. In
particular adding the boundaried-monoid operations introduces no new
collection identities.
\end{corollary}
\begin{proof}
By Proposition~\ref{prop:equiv} it suffices to argue for $\mathbf{T}^{\flat}$,
whose operations are disjoint in sort from those of the comprehension layer and
whose axioms (F1)--(F3) involve no comprehension former; hence they cannot
entail a new identity among comprehension terms.
\end{proof}

%================================================================
\section{Remarks and scope}\label{sec:remarks}

\paragraph{Free models and chains.}
The free boundaried monoid on a family of generating elements with prescribed
boundaries has, as its elements over a pair $(a,b)$, the finite multipliable
chains from $a$ to $b$, with the empty chain at each $a$ as neutral element and
concatenation as product. The passage from generators to chains is itself a
comprehension: a chain is collected by iterating the multipliability guard along
its links, exactly the list-forming instance of the comprehension calculus of
\S\ref{sec:comp}. Thus both the carving of multipliable pairs and the
generation of free elements are handled by the same device.

\paragraph{Determinacy is optional.}
Dropping (T4) yields a strictly weaker theory in which $\kk(g,f)$ may be a
genuine collection rather than a subsingleton: a ``multivalued'' boundaried
monoid, in which several products of a multipliable pair coexist. The remaining
axioms still govern boundaries, support, neutrality, and associativity. We
record this only to note that (T4) is the single axiom responsible for
single-valuedness; it is the natural knob to turn when one later wishes to
relax it.

\paragraph{What the comprehension structure may vary.}
Nothing in \S\ref{sec:T}--\S\ref{sec:equiv} fixes \emph{which} comprehension
structure underlies $\M$, beyond the equations (C1)--(C8). Varying that
structure---while keeping the equations (T1)--(T9) verbatim---yields a family of
theories sharing the present axiomatics. The systematic study of that family,
and of the interpretations it admits, is deferred; here we have only fixed the
equations and shown the partial and total presentations to be one theory.

\paragraph{Summary.}
The theory $\mathbf{T}$ of Definition~\ref{def:axioms} is a purely equational,
finitely axiomatised algebra. Its single design choice is to absorb the
partiality of a boundaried product into a comprehension structure, making the
product a total operation valued in a collection sort. Proposition~\ref{prop:equiv}
shows this total presentation and the dependently sorted partial presentation
to be definitionally equivalent, with comprehension as the explicit translation
between them.

%================================================================
\begin{thebibliography}{9}
\bibitem{wadler}
P.~Wadler.
\emph{Comprehending monads.}
Mathematical Structures in Computer Science \textbf{2}(4):461--493, 1992.
\bibitem{moggi}
E.~Moggi.
\emph{Notions of computation and monads.}
Information and Computation \textbf{93}(1):55--92, 1991.
\end{thebibliography}

\end{document}
